% Options for packages loaded elsewhere
\PassOptionsToPackage{unicode}{hyperref}
\PassOptionsToPackage{hyphens}{url}
\documentclass[
]{article}
\usepackage{xcolor}
\usepackage[margin=1in]{geometry}
\usepackage{amsmath,amssymb}
\setcounter{secnumdepth}{-\maxdimen} % remove section numbering
\usepackage{iftex}
\ifPDFTeX
  \usepackage[T1]{fontenc}
  \usepackage[utf8]{inputenc}
  \usepackage{textcomp} % provide euro and other symbols
\else % if luatex or xetex
  \usepackage{unicode-math} % this also loads fontspec
  \defaultfontfeatures{Scale=MatchLowercase}
  \defaultfontfeatures[\rmfamily]{Ligatures=TeX,Scale=1}
\fi
\usepackage{lmodern}
\ifPDFTeX\else
  % xetex/luatex font selection
\fi
% Use upquote if available, for straight quotes in verbatim environments
\IfFileExists{upquote.sty}{\usepackage{upquote}}{}
\IfFileExists{microtype.sty}{% use microtype if available
  \usepackage[]{microtype}
  \UseMicrotypeSet[protrusion]{basicmath} % disable protrusion for tt fonts
}{}
\makeatletter
\@ifundefined{KOMAClassName}{% if non-KOMA class
  \IfFileExists{parskip.sty}{%
    \usepackage{parskip}
  }{% else
    \setlength{\parindent}{0pt}
    \setlength{\parskip}{6pt plus 2pt minus 1pt}}
}{% if KOMA class
  \KOMAoptions{parskip=half}}
\makeatother
\usepackage{color}
\usepackage{fancyvrb}
\newcommand{\VerbBar}{|}
\newcommand{\VERB}{\Verb[commandchars=\\\{\}]}
\DefineVerbatimEnvironment{Highlighting}{Verbatim}{commandchars=\\\{\}}
% Add ',fontsize=\small' for more characters per line
\usepackage{framed}
\definecolor{shadecolor}{RGB}{248,248,248}
\newenvironment{Shaded}{\begin{snugshade}}{\end{snugshade}}
\newcommand{\AlertTok}[1]{\textcolor[rgb]{0.94,0.16,0.16}{#1}}
\newcommand{\AnnotationTok}[1]{\textcolor[rgb]{0.56,0.35,0.01}{\textbf{\textit{#1}}}}
\newcommand{\AttributeTok}[1]{\textcolor[rgb]{0.13,0.29,0.53}{#1}}
\newcommand{\BaseNTok}[1]{\textcolor[rgb]{0.00,0.00,0.81}{#1}}
\newcommand{\BuiltInTok}[1]{#1}
\newcommand{\CharTok}[1]{\textcolor[rgb]{0.31,0.60,0.02}{#1}}
\newcommand{\CommentTok}[1]{\textcolor[rgb]{0.56,0.35,0.01}{\textit{#1}}}
\newcommand{\CommentVarTok}[1]{\textcolor[rgb]{0.56,0.35,0.01}{\textbf{\textit{#1}}}}
\newcommand{\ConstantTok}[1]{\textcolor[rgb]{0.56,0.35,0.01}{#1}}
\newcommand{\ControlFlowTok}[1]{\textcolor[rgb]{0.13,0.29,0.53}{\textbf{#1}}}
\newcommand{\DataTypeTok}[1]{\textcolor[rgb]{0.13,0.29,0.53}{#1}}
\newcommand{\DecValTok}[1]{\textcolor[rgb]{0.00,0.00,0.81}{#1}}
\newcommand{\DocumentationTok}[1]{\textcolor[rgb]{0.56,0.35,0.01}{\textbf{\textit{#1}}}}
\newcommand{\ErrorTok}[1]{\textcolor[rgb]{0.64,0.00,0.00}{\textbf{#1}}}
\newcommand{\ExtensionTok}[1]{#1}
\newcommand{\FloatTok}[1]{\textcolor[rgb]{0.00,0.00,0.81}{#1}}
\newcommand{\FunctionTok}[1]{\textcolor[rgb]{0.13,0.29,0.53}{\textbf{#1}}}
\newcommand{\ImportTok}[1]{#1}
\newcommand{\InformationTok}[1]{\textcolor[rgb]{0.56,0.35,0.01}{\textbf{\textit{#1}}}}
\newcommand{\KeywordTok}[1]{\textcolor[rgb]{0.13,0.29,0.53}{\textbf{#1}}}
\newcommand{\NormalTok}[1]{#1}
\newcommand{\OperatorTok}[1]{\textcolor[rgb]{0.81,0.36,0.00}{\textbf{#1}}}
\newcommand{\OtherTok}[1]{\textcolor[rgb]{0.56,0.35,0.01}{#1}}
\newcommand{\PreprocessorTok}[1]{\textcolor[rgb]{0.56,0.35,0.01}{\textit{#1}}}
\newcommand{\RegionMarkerTok}[1]{#1}
\newcommand{\SpecialCharTok}[1]{\textcolor[rgb]{0.81,0.36,0.00}{\textbf{#1}}}
\newcommand{\SpecialStringTok}[1]{\textcolor[rgb]{0.31,0.60,0.02}{#1}}
\newcommand{\StringTok}[1]{\textcolor[rgb]{0.31,0.60,0.02}{#1}}
\newcommand{\VariableTok}[1]{\textcolor[rgb]{0.00,0.00,0.00}{#1}}
\newcommand{\VerbatimStringTok}[1]{\textcolor[rgb]{0.31,0.60,0.02}{#1}}
\newcommand{\WarningTok}[1]{\textcolor[rgb]{0.56,0.35,0.01}{\textbf{\textit{#1}}}}
\usepackage{longtable,booktabs,array}
\usepackage{calc} % for calculating minipage widths
% Correct order of tables after \paragraph or \subparagraph
\usepackage{etoolbox}
\makeatletter
\patchcmd\longtable{\par}{\if@noskipsec\mbox{}\fi\par}{}{}
\makeatother
% Allow footnotes in longtable head/foot
\IfFileExists{footnotehyper.sty}{\usepackage{footnotehyper}}{\usepackage{footnote}}
\makesavenoteenv{longtable}
\usepackage{graphicx}
\makeatletter
\newsavebox\pandoc@box
\newcommand*\pandocbounded[1]{% scales image to fit in text height/width
  \sbox\pandoc@box{#1}%
  \Gscale@div\@tempa{\textheight}{\dimexpr\ht\pandoc@box+\dp\pandoc@box\relax}%
  \Gscale@div\@tempb{\linewidth}{\wd\pandoc@box}%
  \ifdim\@tempb\p@<\@tempa\p@\let\@tempa\@tempb\fi% select the smaller of both
  \ifdim\@tempa\p@<\p@\scalebox{\@tempa}{\usebox\pandoc@box}%
  \else\usebox{\pandoc@box}%
  \fi%
}
% Set default figure placement to htbp
\def\fps@figure{htbp}
\makeatother
\setlength{\emergencystretch}{3em} % prevent overfull lines
\providecommand{\tightlist}{%
  \setlength{\itemsep}{0pt}\setlength{\parskip}{0pt}}
\usepackage{bookmark}
\IfFileExists{xurl.sty}{\usepackage{xurl}}{} % add URL line breaks if available
\urlstyle{same}
\hypersetup{
  hidelinks,
  pdfcreator={LaTeX via pandoc}}

\author{}
\date{\vspace{-2.5em}}

\begin{document}

\section{EU Forest Observatory: Forest Loss
Monitoring}\label{eu-forest-observatory-forest-loss-monitoring}

Google Earth Engine and R workflow for producing annual forest-loss
layers, aggregating them to a 0.2-degree grid, detecting extreme loss
events, and generating country-level statistics by disturbance type.

\subsection{Overview}\label{overview}

This repository documents a processing chain developed for the EU Forest
Observatory. The workflow combines:

\begin{enumerate}
\def\labelenumi{\arabic{enumi}.}
\tightlist
\item
  Hansen Global Forest Change data for tree cover and annual forest
  loss.
\item
  An annual fire-related forest-loss product used to exclude
  fire-affected pixels from the harvest-oriented layers.
\item
  Google Earth Engine aggregation from approximately 30 m to a nominal
  0.02-degree grid (approximately 2 km).
\item
  A second Earth Engine aggregation from the approximately 2 km layers
  to a 0.2-degree grid (approximately 20 km).
\item
  An R-based robust outlier procedure based on the median and median
  absolute deviation (MAD).
\item
  A final Earth Engine workflow producing country-level annual
  statistics for total forest loss, fire-related loss, and extreme-loss
  events.
\item
  The Curtis et al.~forest-loss-driver map for contextual interpretation
  of disturbance causes.
\end{enumerate}

The analysis follows the conceptual approach described in Ceccherini et
al.~(2020), including spatial aggregation and the separation of abrupt
or extreme disturbances from the normal loss signal. The original Nature
study and its reproducibility materials are available through Zenodo:
\href{https://doi.org/10.5281/zenodo.3687096}{code} and
\href{https://doi.org/10.5281/zenodo.3687090}{data}.

\subsection{Repository structure}\label{repository-structure}

\begin{Shaded}
\begin{Highlighting}[]
\NormalTok{.}
\NormalTok{├── README.md}
\NormalTok{├── gee/}
\NormalTok{│   ├── 01\_prepare\_annual\_loss\_assets.js}
\NormalTok{│   ├── 02\_aggregate\_to\_20km.js}
\NormalTok{│   └── 04\_country\_statistics.js}
\NormalTok{├── R/}
\NormalTok{│   └── 03\_detect\_extreme\_loss.R}
\NormalTok{├── data/}
\NormalTok{│   ├── raw/                 \# Local GeoTIFFs downloaded from Earth Engine}
\NormalTok{│   ├── intermediate/        \# R{-}derived rasters and intermediate products}
\NormalTok{│   └── README.md}
\NormalTok{├── docs/}
\NormalTok{│   ├── workflow.md}
\NormalTok{│   └── data\_dictionary.md}
\NormalTok{└── CITATION.cff}
\end{Highlighting}
\end{Shaded}

The code supplied for this documentation should be split into the files
above. Asset IDs, Drive folders, regions, country lists, and local paths
must be adapted to the execution account and project.

\subsection{Workflow}\label{workflow}

\subsubsection{1. Prepare annual loss
assets}\label{prepare-annual-loss-assets}

\texttt{gee/01\_prepare\_annual\_loss\_assets.js} loads:

\begin{itemize}
\tightlist
\item
  \texttt{UMD/hansen/global\_forest\_change\_2025\_v1\_13};
\item
  \texttt{users/sashatyu/2001-2025\_fire\_forest\_loss\_annual};
\item
  an analysis region, currently represented by \texttt{AOItot} in the
  script.
\end{itemize}

The Hansen product is a 30 m Landsat-derived dataset covering
2000--2025. Its \texttt{treecover2000} band represents canopy cover in
2000, while \texttt{lossyear} encodes loss years as 1--25 for
2001--2025. The official Earth Engine catalogue documents the product
and its bands at
\href{https://developers.google.com/earth-engine/datasets/catalog/UMD_hansen_global_forest_change_2025_v1_13}{Hansen
Global Forest Change v1.13}.

The script applies the following masks:

\begin{itemize}
\tightlist
\item
  \texttt{treecover2000\ \textgreater{}=\ forest\_threshold}, with the
  default threshold set to 10\%.
\item
  \texttt{gain\ \textless{}\ 1}, excluding pixels classified as forest
  gain.
\item
  Annual fire mask equal to zero, excluding pixels identified by the
  fire product.
\end{itemize}

For each year, a binary loss layer is produced using
\texttt{lossyear.eq(year\_code)}. The script also constructs a
cumulative forest-presence sequence beginning in 2000. These
forest-presence layers are intended to represent the initial forest
denominator for subsequent aggregation.

Each binary annual loss layer and the 2000 forest layer is aggregated
with:

\begin{Shaded}
\begin{Highlighting}[]
\OperatorTok{.}\FunctionTok{reduceResolution}\NormalTok{(ee}\OperatorTok{.}\AttributeTok{Reducer}\OperatorTok{.}\FunctionTok{sum}\NormalTok{()}\OperatorTok{.}\FunctionTok{unweighted}\NormalTok{()}\OperatorTok{,} \KeywordTok{false}\OperatorTok{,} \DecValTok{65536}\NormalTok{)}
\OperatorTok{.}\FunctionTok{reproject}\NormalTok{(ee}\OperatorTok{.}\FunctionTok{Projection}\NormalTok{(}\StringTok{\textquotesingle{}EPSG:4326\textquotesingle{}}\NormalTok{)}\OperatorTok{.}\FunctionTok{scale}\NormalTok{(}\FloatTok{0.02}\OperatorTok{,} \FloatTok{0.02}\NormalTok{))}
\end{Highlighting}
\end{Shaded}

The resulting values are sums of source pixels within the target grid
cell. They should therefore be interpreted as counts of approximately 30
m pixels, not directly as hectares or percentages.

The script exports one Earth Engine asset per year and one 2000
forest-denominator asset. The export names follow this pattern:

\begin{Shaded}
\begin{Highlighting}[]
\NormalTok{Forest2000\_at\_2km\_2025GlobalFires\_10}
\NormalTok{loss\_YYYY\_at\_2km\_2025GlobalFires\_10}
\end{Highlighting}
\end{Shaded}

Despite the historical variable names \texttt{at\_2km}, the target
projection uses 0.02 degrees. At the equator this is approximately 2.2
km, while the physical north--south and east--west dimensions vary with
latitude.

\subsubsection{2. Aggregate annual assets to approximately 20
km}\label{aggregate-annual-assets-to-approximately-20-km}

\texttt{gee/02\_aggregate\_to\_20km.js} loads the annual approximately 2
km loss assets, stacks them into \texttt{Final\_loss}, and applies a
second sum aggregation to a 0.2-degree grid:

\begin{Shaded}
\begin{Highlighting}[]
\NormalTok{Final\_loss}
  \OperatorTok{.}\FunctionTok{reduceResolution}\NormalTok{(ee}\OperatorTok{.}\AttributeTok{Reducer}\OperatorTok{.}\FunctionTok{sum}\NormalTok{()}\OperatorTok{.}\FunctionTok{unweighted}\NormalTok{()}\OperatorTok{,} \KeywordTok{false}\OperatorTok{,} \DecValTok{65536}\NormalTok{)}
  \OperatorTok{.}\FunctionTok{reproject}\NormalTok{(ee}\OperatorTok{.}\FunctionTok{Projection}\NormalTok{(}\StringTok{\textquotesingle{}EPSG:4326\textquotesingle{}}\NormalTok{)}\OperatorTok{.}\FunctionTok{scale}\NormalTok{(}\FloatTok{0.2}\OperatorTok{,} \FloatTok{0.2}\NormalTok{))}
\end{Highlighting}
\end{Shaded}

The script exports:

\begin{itemize}
\tightlist
\item
  \texttt{FinalLoss\_at\_20km\_2025Fires.tif}, containing one band per
  annual loss layer;
\item
  \texttt{Forest2000\_at\_20km\_2025Fires.tif}, containing the
  aggregated 2000 forest denominator.
\end{itemize}

This is a nominal scale corresponding to 0.2 degrees at the equator; it
is not a constant metric 20 km grid globally.

The current supplied script loads loss assets from 2004 onward.

\subsubsection{3. Detect extreme loss events in
R}\label{detect-extreme-loss-events-in-r}

\texttt{R/03\_detect\_extreme\_loss.R} reads the two 0.2-degree GeoTIFFs
with the \texttt{raster} package:

\begin{Shaded}
\begin{Highlighting}[]
\NormalTok{Final\_loss }\OtherTok{\textless{}{-}} \FunctionTok{stack}\NormalTok{(}\StringTok{"Data2025/FinalLoss\_at\_20km\_2025Fires.tif"}\NormalTok{)}
\NormalTok{Forest }\OtherTok{\textless{}{-}} \FunctionTok{stack}\NormalTok{(}\StringTok{"Data2025/Forest2000\_at\_20km\_2025Fires.tif"}\NormalTok{)}
\end{Highlighting}
\end{Shaded}

The script first calculates the annual loss percentage relative to the
2000 forest-pixel denominator:

\begin{Shaded}
\begin{Highlighting}[]
\NormalTok{Rho }\OtherTok{\textless{}{-}}\NormalTok{ (Final\_loss }\SpecialCharTok{/}\NormalTok{ Forest[[}\DecValTok{1}\NormalTok{]]) }\SpecialCharTok{*} \DecValTok{100}
\NormalTok{Rho[Rho }\SpecialCharTok{==} \DecValTok{0}\NormalTok{] }\OtherTok{\textless{}{-}} \ConstantTok{NA}
\end{Highlighting}
\end{Shaded}

It then calculates cell-wise temporal summary statistics:

\begin{itemize}
\tightlist
\item
  median;
\item
  mean;
\item
  MAD;
\item
  standard deviation.
\end{itemize}

The classification rule is:

\begin{Shaded}
\begin{Highlighting}[]
\NormalTok{Rho }\SpecialCharTok{\textgreater{}}\NormalTok{ RhoM }\SpecialCharTok{+} \DecValTok{3} \SpecialCharTok{*}\NormalTok{ Rhomad }\SpecialCharTok{\&}
\NormalTok{Rho }\SpecialCharTok{\textgreater{}} \DecValTok{3} \SpecialCharTok{\&}
\NormalTok{(Forest }\SpecialCharTok{/} \DecValTok{640000}\NormalTok{) }\SpecialCharTok{\textgreater{}} \FloatTok{0.05}
\end{Highlighting}
\end{Shaded}

A cell is classified as an extreme-loss event when its annual relative
loss exceeds the temporal median by three MADs, exceeds 3\%, and
contains more than 5\% forest cover according to the approximate
source-pixel denominator. The binary mask is written to:

\begin{Shaded}
\begin{Highlighting}[]
\NormalTok{Data2025/MASKGEE2025Fires.tif}
\end{Highlighting}
\end{Shaded}

The continuous relative-loss time series is written to:

\begin{Shaded}
\begin{Highlighting}[]
\NormalTok{Data2025/ForestHarvest\_04\_25Fires.tif}
\end{Highlighting}
\end{Shaded}

The R-derived mask is then uploaded to Earth Engine as the asset used by
the country-statistics script. Please note that the time series starts
in 2004.

\subsubsection{4. Produce country-level
statistics}\label{produce-country-level-statistics}

\texttt{gee/04\_country\_statistics.js} combines:

\begin{itemize}
\tightlist
\item
  the R-derived extreme-event mask;
\item
  the Curtis driver map;
\item
  the Hansen tree-cover and loss bands;
\item
  annual fire-loss layers;
\item
  country boundaries;
\item
  country-specific tree-cover thresholds.
\end{itemize}

The Curtis layer is loaded as:

\begin{Shaded}
\begin{Highlighting}[]
\NormalTok{ee}\OperatorTok{.}\FunctionTok{Image}\NormalTok{(}\StringTok{"projects/tmf{-}monitoring/assets/CurtisDrivers2018/FilledMap"}\NormalTok{)}
\end{Highlighting}
\end{Shaded}

The code restricts the Hansen analysis to \texttt{CURTIS.eq(3)}, which
refers to ``forestry''. The script then builds a year-coded
extreme-event image from the multiband R mask. Bands \texttt{b8} through
\texttt{b22} are mapped to years 2011--2025 using the values 11--25.

For each country, the script applies the corresponding country-specific
threshold from \texttt{list\_t}. It then calculates annual area
statistics at 30 m using pixel area and grouped reducers.

Three CSV outputs are produced:

\begin{enumerate}
\def\labelenumi{\arabic{enumi}.}
\tightlist
\item
  \texttt{Country\_Forest\_Change\_EUOBS\_\textless{}country\textgreater{}.csv}
  --- all forest loss by loss year.
\item
  \texttt{Country\_Forest\_Change\_EUOBS\_fires\textless{}country\textgreater{}.csv}
  --- fire-related loss, excluding cells classified as extreme events by
  the wind/extreme mask.
\item
  \texttt{Country\_Forest\_Change\_EUOBS\_Wind\textless{}country\textgreater{}.csv}
  --- loss in cells classified as extreme events (i.e.~Wind and other
  extreme events).
\end{enumerate}

The grouped reducer converts the \texttt{lossyear} or annual
fire/extreme-event code into wide columns such as \texttt{sum\_1},
\texttt{sum\_2}, and so forth. The values are areas in square metres
because loss pixels are multiplied by \texttt{ee.Image.pixelArea()}.

\subsubsection{Fire and extreme-event
terminology}\label{fire-and-extreme-event-terminology}

The supplied variable names use \texttt{WIND}, but the R mask is
generated from temporal outliers in the aggregated loss series. It
should therefore be called an \texttt{extreme\_loss} or
\texttt{abrupt\_loss} mask unless it has been independently validated as
wind damage. If the mask is specifically intended to represent
windstorms, document the external windstorm layer, temporal coverage,
and exclusion rule.

\subsection{Reproducible execution}\label{reproducible-execution}

\begin{enumerate}
\def\labelenumi{\arabic{enumi}.}
\tightlist
\item
  Open \texttt{gee/01\_prepare\_annual\_loss\_assets.js} in the Earth
  Engine Code Editor.
\item
  Define \texttt{AOItot}, confirm the Hansen and fire asset versions,
  and start the annual asset exports.
\item
  Wait until the asset exports are complete and confirm their band names
  and footprints.
\item
  Update the asset IDs in \texttt{gee/02\_aggregate\_to\_20km.js} and
  export the two GeoTIFFs to Google Drive.
\item
  Download the GeoTIFFs into the local \texttt{data/raw/} directory.
\item
  Run \texttt{R/03\_detect\_extreme\_loss.R} in RStudio and inspect the
  denominator, relative-loss, MAD, and extreme-event maps.
\item
  Upload \texttt{MASKGEE2025Fires.tif} to Earth Engine and update its
  project asset ID in \texttt{gee/04\_country\_statistics.js}.
\item
  Define the country boundary collection and verify the country-code
  vectors and threshold vector have identical lengths and ordering.
\item
  Run the country-statistics script and retrieve the generated CSV files
  from Google Drive.
\item
  Archive the Earth Engine task configuration, asset IDs, input
  versions, thresholds, and output checksums.
\end{enumerate}

\subsection{Data dictionary}\label{data-dictionary}

\begin{longtable}[]{@{}
  >{\raggedright\arraybackslash}p{(\linewidth - 4\tabcolsep) * \real{0.3333}}
  >{\raggedright\arraybackslash}p{(\linewidth - 4\tabcolsep) * \real{0.3333}}
  >{\raggedright\arraybackslash}p{(\linewidth - 4\tabcolsep) * \real{0.3333}}@{}}
\toprule\noalign{}
\begin{minipage}[b]{\linewidth}\raggedright
Object
\end{minipage} & \begin{minipage}[b]{\linewidth}\raggedright
Type
\end{minipage} & \begin{minipage}[b]{\linewidth}\raggedright
Meaning
\end{minipage} \\
\midrule\noalign{}
\endhead
\bottomrule\noalign{}
\endlastfoot
\texttt{treecover2000} & Hansen band & Tree canopy cover in 2000,
percent. \\
\texttt{gain} & Hansen band & Forest gain flag for 2000--2012; not
updated in later versions. \\
\texttt{lossyear} & Hansen band & Loss-year code 1--25 for
2001--2025. \\
\texttt{forest\_threshold} & Parameter & Tree-cover threshold used to
define forest. \\
\texttt{Final\_loss} & Earth Engine image & Annual aggregated loss bands
at the intermediate grid. \\
\texttt{Forest} & Earth Engine image & Aggregated 2000 forest-pixel
denominator. \\
\texttt{Rho} & R raster stack & Annual loss as a percentage of the 2000
forest denominator. \\
\texttt{RhoM} & R raster & Temporal median of \texttt{Rho} per grid
cell. \\
\texttt{Rhomad} & R raster & Temporal MAD of \texttt{Rho} per grid
cell. \\
\texttt{Rho2} & R raster & Binary extreme-loss mask. \\
\texttt{WIND\_TOT} & Earth Engine image & Year-coded extreme-event layer
reconstructed from the uploaded mask. \\
\texttt{CURTIS} & Earth Engine image & Curtis forest-loss-driver map. \\
\texttt{pixelArea()} & Earth Engine image & Pixel area used for
country-level area totals. \\
\end{longtable}

\subsection{Scientific context}\label{scientific-context}

The Hansen GFC product is a Landsat-based global forest-change dataset
with approximately 30 m pixels and annual loss-year coding through 2025
in the current v1.13 release. The product is documented in the Earth
Engine Data Catalog and is distributed under CC BY 4.0. The original
methodological basis is Hansen et al.~(2013).

Curtis et al.~(2018) classified dominant drivers of global forest loss,
including commodity-driven deforestation, shifting agriculture,
forestry, wildfire, and urbanisation. The driver layer in this
repository is used as a spatial stratification or filter, not as a
substitute for an independent validation of each detected loss event.

The Nature study by Ceccherini et al.~(2020) used aggregated
satellite-derived forest-loss information to study harvested forest area
in Europe, while excluding fire-affected areas and discussing the
treatment of major windstorms. This repository extends that general
structure with the 2025 Hansen release and a robust time-series
procedure for identifying unusually large grid-cell losses.

\subsection{Citation}\label{citation}

\begin{itemize}
\tightlist
\item
  Ceccherini, G., Duveiller, G., Grassi, G., et al.~(2020). Abrupt
  increase in harvested forest area over Europe after 2015.
  \emph{Nature}, 583, 72--77.
  \url{https://doi.org/10.1038/s41586-020-2438-y}
\item
  Curtis, P. G., Slay, C. M., Harris, N. L., Tyukavina, A., and Hansen,
  M. C. (2018). Classifying drivers of global forest loss.
  \emph{Science}, 361, 1108--1111.
  \url{https://doi.org/10.1126/science.aau3445}
\item
  Hansen, M. C., et al.~(2013). High-resolution global maps of
  21st-century forest cover change. \emph{Science}, 342, 850--853.
  \url{https://doi.org/10.1126/science.1244693}
\end{itemize}

\end{document}
